\BeleskeID{09_3-s050}
% element: 09_3-s050:e04
% slide-id: 09_3-s050
Smanjenje broja tranzistora naročito je moguće kod selektorskih funkcija: za jednu vrednost upravljačkog signala prosleđuje se jedan signal, a za komplementarnu drugi. Osnovna ideja prolazne logike je da tranzistor radi kao prekidač u putanji signala.

Izlaz se podrazumeva kao kapacitivno opterećen. Zbog čuvanja naelektrisanja prolazna kola često se svrstavaju u dinamička, iako je njihov princip prenosa drugačiji od pripreme i izračunavanja u dinamičkoj logici.
